\chapter{Risultati}
\label{cap:risultati}

Tutti i numeri di questo capitolo seguono il protocollo fissato nel
Capitolo~\ref{cap:metriche}: $2\,000$ prompt del test set, $N=5$ generazioni per
prompt, campionamento a temperatura $0{,}7$, metriche versione $2$.

\section{Statuto dei risultati: valori provvisori}
\label{sec:statuto-provvisorio}

Questo avvertimento precede la tabella perché ne condiziona la lettura, e
ometterlo renderebbe il capitolo fuorviante.

\textbf{I valori riportati provengono da esecuzioni che precedono le revisioni
del codice descritte in questa relazione.} Nell'intervallo fra la loro
produzione e la stesura di questo testo la pipeline è cambiata in modi che
influiscono su ciò che si misura, non soltanto su come lo si misura:

\begin{center}
\begin{tabular}{lll}
\toprule
Parametro & Quando i numeri furono prodotti & Configurazione corrente \\
\midrule
Passi di ottimizzazione (GRPO) & $2\,000$ & $5\,000$ \\
Prompt valutati & $2\,000$ & $5\,000$ \\
Modalità di valutazione & singola & doppia (entrambi i prompting) \\
Curriculum per difficoltà & attivo sulle celle few-shot & rimosso \\
Retrieval few-shot & implementazione originale & riscritta (equivalente) \\
\bottomrule
\end{tabular}
\end{center}

Le prime tre voci cambiano il numero. Passare da $2\,000$ a $5\,000$ passi
significa che il modello vede $40\,000$ prompt invece di $16\,000$, cioè il
$54{,}8\%$ del training set invece del $21{,}9\%$: le celle addestrate con
reinforcement learning \textbf{non sono le stesse}. Passare a $5\,000$ prompt di
valutazione restringe gli intervalli di confidenza ma rende i valori non
direttamente confrontabili con quelli qui riportati.

La quarta voce, la rimozione del curriculum, non dovrebbe cambiare nulla,
perché quel componente era inerte; ma la sua inerzia è stata stabilita dopo che
questi numeri erano già stati prodotti.

La quinta è l'unica per cui esiste una garanzia forte: la riscrittura del
retrieval è stata verificata equivalente elemento per elemento su migliaia di
interrogazioni, quindi non altera i risultati.

\subsection{Che cosa resta valido}

La conclusione centrale della relazione \textbf{non} dipende da questi numeri
nel dettaglio, e questo è il motivo per cui vale la pena pubblicarli comunque.
Le affermazioni che sopravvivono a una loro revisione sono:

\begin{itemize}
  \item la baseline a regole ottiene ROUGE-L $\approx 0{,}97$, ed è un
        riferimento indipendente da qualunque addestramento
        (sezione~\ref{sec:baseline-regole});
  \item nessuna cella con reinforcement learning si avvicina a quel valore, e
        il divario è di ordine $0{,}35$, cioè cento volte l'incertezza
        sperimentale;
  \item la trasformazione banale $\mathrm{uppercase}(\text{text})$ ottiene
        $0{,}6454$, e per superare quel valore un metodo di rinforzo dovrebbe
        migliorare di una quantità molto maggiore di qualunque variazione
        plausibile fra run;
  \item il meccanismo del risultato negativo --- il collasso del supporto,
        misurato sui rollout (sezione~\ref{sec:supporto}) --- è una proprietà
        del substrato e della reward, non di un particolare numero di passi.
\end{itemize}

Va invece considerato \textbf{provvisorio} tutto ciò che dipende da differenze
piccole: il confronto fra celle vicine (per esempio le ablazioni, che
differiscono di qualche millesimo), e ogni valore riportato alla quarta cifra
decimale. La sezione~\ref{sec:dispersione-run} mostra del resto che la
dispersione fra due esecuzioni della stessa cella può superare l'intervallo di
confidenza: la precisione apparente di queste tabelle è già oggi superiore alla
loro accuratezza.

\section{Provenienza dei numeri}
\label{sec:provenienza}

Stabilito lo statuto, va dichiarata la provenienza, perché una precedente
stesura di questo capitolo mescolava versioni diverse del codice di valutazione
e va evitato che l'errore si ripeta.

Ogni riga della tabella principale è letta da \emph{un solo} file
\texttt{eval\_*.json} prodotto dalla pipeline di valutazione, e per ogni cella è
usato il run più recente disponibile. Tutti i file dichiarano
\texttt{metrics\_version = 2}. Nessun valore proviene da documenti derivati o da
ri-calcoli successivi.

Questa precisazione è necessaria perché nel progetto è esistita una versione $3$
sperimentale del modulo delle metriche, con definizioni diverse di validity e di
exact match normalizzato. La versione $3$ \textbf{non} è uniformemente più
severa della $2$: sull'uscita non vincolata è più permissiva (validity
$0{,}0370$ in v2 contro $0{,}1043$ in v3), mentre sull'uscita quasi corretta
dell'SFT è più severa sull'exact match ($0{,}8117$ in v2 contro un valore
inferiore in v3, per effetto della normalizzazione degli spazi). Mescolare le due
versioni nella stessa tabella produce quindi incoerenze in entrambe le
direzioni, ed è esattamente ciò che va evitato.

Una nota separata riguarda la \emph{non-copy-token accuracy}. Il
Capitolo~\ref{cap:metriche} la promuove a metrica primaria, ma nelle esecuzioni
da cui provengono queste tabelle essa \textbf{non} era calcolata dalla pipeline
di valutazione: la sua implementazione viveva in un percorso di analisi
separato, eseguito a posteriori sulle generazioni salvate. I valori riportati
provengono da lì. La metrica è stata in seguito integrata nella pipeline, quindi
le esecuzioni future la produrranno insieme alle altre; per le celle di queste
tabelle resta un valore ricostruito.

\section{La tabella principale}

\begin{center}
\begin{tabular}{lrrrrrr}
\toprule
Cella & ROUGE-L & Pass@1 & BLEU-c & GlossF1 & Validity & EM \\
\midrule
Base zero-shot senza Trie & $0{,}3648$ & $0{,}5845$ & $0{,}0008$ & $0{,}2279$ & $0{,}0370$ & $0{,}0007$ \\
Base zero-shot con Trie & $0{,}1373$ & $0{,}1820$ & $0{,}0270$ & $0{,}1186$ & $0{,}9055$ & $0{,}0017$ \\
Base few-shot con Trie & $0{,}4664$ & $0{,}7330$ & $0{,}1880$ & $0{,}4026$ & $0{,}9821$ & $0{,}0148$ \\
GRPO da modello base & $0{,}5998$ & $0{,}9315$ & $0{,}3263$ & $0{,}6072$ & $0{,}9933$ & $0{,}0499$ \\
SFT seguito da GRPO & $0{,}5114$ & $0{,}8100$ & $0{,}2222$ & $0{,}4702$ & $0{,}9925$ & $0{,}0203$ \\
SFT soltanto & $0{,}9749$ & $0{,}9980$ & $0{,}9426$ & $0{,}9760$ & $0{,}9998$ & $0{,}8117$ \\
\midrule
\textbf{Baseline a regole} & $\mathbf{0{,}9685}$ & --- & $\mathbf{0{,}8907}$ & $\mathbf{0{,}9665}$ & --- & $\mathbf{0{,}5865}$ \\
$\mathrm{uppercase}(\text{text})$ & $0{,}6454$ & --- & --- & --- & --- & --- \\
\bottomrule
\end{tabular}
\end{center}

La riga in grassetto è il riferimento che rende il resto leggibile. Le
osservazioni sostanziali sono cinque, ed è opportuno affrontarle in ordine di
importanza decrescente.

Due avvertenze sulla lettura della tabella.

La colonna BLEU riporta il valore \emph{a corpus}. Va notato che i valori di
BLEU e chrF sono calcolati su dato pre-tokenizzato: il gloss ha la punteggiatura
come token separato (\texttt{X-Y WILL DESC-NOT FLINCH .}), e la libreria
impiegata lo segnala. I confronti \emph{fra} celle di questa tabella restano
validi, perché il protocollo è identico per tutte; i valori assoluti
\textbf{non} sono confrontabili con numeri pubblicati calcolati su dato
detokenizzato.

La colonna Pass@1 usa una soglia permissiva (ROUGE-L $\ge 0{,}3$, si veda il
Capitolo~\ref{cap:metriche}). Il primo valore della colonna, $0{,}5845$ sulla
cella senza vincolo, va letto insieme all'exact match della stessa riga
($0{,}0007$) e alla sua validity ($0{,}0370$): una cella in cui il $96\%$ delle
uscite non è nemmeno una sequenza di gloss ben formata ottiene comunque
pass@1 $0{,}58$. È una misura di copertura del campionamento rispetto a una
soglia bassa, non di qualità.

\section{Osservazione 1: nessun risultato con RL supera la regola}

Il miglior risultato ottenuto con reinforcement learning è $0{,}5998$ di ROUGE-L
(GRPO dal modello base). La baseline a regole ottiene $0{,}9685$. La
trasformazione \emph{metti tutto in maiuscolo} ottiene $0{,}6454$.

Cioè: il miglior risultato con RL è \textbf{inferiore} a quello ottenibile
maiuscolizzando l'input, e distante $0{,}36$ da una tabella di corrispondenze
costruita contando parole nel training set.

Su exact match il divario è più netto: $0{,}0499$ contro $0{,}5865$, un fattore
dodici. Su gloss F1: $0{,}6072$ contro $0{,}9665$.

Non esiste una lettura in cui questi numeri rappresentino un successo del
reinforcement learning su questo compito. La domanda residua non è ``quanto
migliora'', ma ``perché non poteva funzionare''; il capitolo sul GRPO ha dato la
risposta meccanica, e la sezione~\ref{sec:decomposizione} la quantifica.

\section{Osservazione 2: l'SFT è al tetto, il che invalida la metrica}

L'SFT ottiene ROUGE-L $0{,}9749$, superiore sia alla baseline a regole
($0{,}9685$) sia al miglior valore pubblicato in letteratura su ASLG-PC12
($0{,}9587$, \cite{peng2023monolingual}).

Un modello da $0{,}5$ miliardi di parametri, con adattatori LoRA, addestrato per
tre epoche su una singola GPU, che supera lo stato dell'arte pubblicato. La
lettura corretta non è che il nostro sistema sia migliore: è che la metrica ha
esaurito la propria capacità discriminante.

Il contrasto con PHOENIX-2014T rende la cosa evidente. Su quel corpus, dove la
sovrapposizione lessicale fra gloss e testo è $15{,}59\%$ invece di $98{,}23\%$
(sezione~\ref{sec:sintetico}), i migliori risultati T2G riportati sono intorno a
ROUGE-L $55{,}18$~\cite{ouargani2023t2g}. Circa quaranta punti in meno, sullo
stesso compito nominale e con metodi comparabili. La differenza non è nei
metodi: è nel dato.

\section{Osservazione 3: il vincolo simbolico e le due metriche}

Le prime due righe della tabella principale, lette insieme alla non-copy-token
accuracy, sono il risultato neuro-simbolico del lavoro.

\begin{center}
\begin{tabular}{lrrr}
\toprule
Cella & ROUGE-L & Validity & Non-copy accuracy \\
\midrule
Base zero-shot senza Trie & $0{,}3648$ & $0{,}0370$ & $0{,}0006$ \\
Base zero-shot con Trie & $0{,}1373$ & $0{,}9055$ & $0{,}0432$ \\
\bottomrule
\end{tabular}
\end{center}

Il vincolo abbassa ROUGE-L di $0{,}23$, alza la validity di un fattore $24$
(da $0{,}0370$ a $0{,}9055$) e alza la non-copy accuracy di un fattore circa
$70$. Il meccanismo è quello discusso nel Capitolo~\ref{cap:decoding}: senza
vincolo il modello produce inglese in minuscolo, che la metrica premia per il
difetto documentato nella sezione~\ref{sec:difetto-rouge}; con il vincolo
produce gloss, male, ma producendo gloss.

La validity di $0{,}0370$ nella prima riga è il dato che chiude la questione:
il $96\%$ delle uscite del modello senza vincolo non è nemmeno una sequenza di
gloss ben formata. Assegnare loro ROUGE-L $0{,}3648$ è un artefatto della
metrica, non una misura di qualità.

\subsection{Una conferma indipendente dalla scelta della metrica}
\label{sec:conferma-case-sensitive}

Su questa cella esiste una verifica che non dipende dalla non-copy accuracy né
da alcuna metrica introdotta da questo lavoro, e che è quindi la prova più
difficile da contestare. Sulle \emph{stesse} generazioni, senza vincolo:

\begin{center}
\begin{tabular}{lrl}
\toprule
Metrica & Valore & Sensibile al maiuscolo? \\
\midrule
ROUGE-L & $0{,}3648$ & no \\
chrF (a corpus) & $1{,}18$ su $100$ & sì \\
BLEU (a corpus) & $0{,}0008$ & sì \\
\bottomrule
\end{tabular}
\end{center}

Le tre metriche sono calcolate sullo stesso file di generazioni, nella stessa
esecuzione. Le due sensibili al maiuscolo assegnano un valore prossimo a zero;
quella insensibile assegna $0{,}36$. Il rapporto è di circa un fattore $30$.

Questo è il difetto della sezione~\ref{sec:difetto-rouge} osservato su
$10\,000$ generazioni reali invece che su casi costruiti a mano: il modello
produce inglese minuscolo, ROUGE-L lo premia perché normalizza via il maiuscolo,
e ogni metrica che il maiuscolo lo guarda registra il fallimento. Per confronto,
sulla cella con vincolo attivo chrF sale a $13{,}21$, cioè di un fattore $11$
rispetto a $1{,}18$ --- nella direzione opposta a quella indicata da ROUGE-L.

\section{Osservazione 4: le lunghezze non mostrano deriva verso la verbosità}

Un'ipotesi corrente sul reinforcement learning è che tenda a produrre uscite più
lunghe, perché la reward correla con la lunghezza~\cite{singhal2023length}. Era
l'ipotesi di partenza dell'analisi della reward
(sezione sul test avversario nel Capitolo~\ref{cap:reward}), e va verificata sui
dati.

\begin{center}
\begin{tabular}{lrrrrr}
\toprule
Cella & Gen.\ & Rif.\ & Rapporto & Errore firmato & Errore assoluto \\
\midrule
GRPO da base & $12{,}40$ & $12{,}26$ & $1{,}011$ & $+0{,}14$ & $2{,}81$ \\
SFT seguito da GRPO & $10{,}86$ & $12{,}26$ & $0{,}886$ & $-1{,}40$ & $3{,}59$ \\
SFT soltanto & $12{,}32$ & $12{,}26$ & $1{,}005$ & $+0{,}06$ & $0{,}09$ \\
Base zero-shot con Trie & $13{,}96$ & $12{,}26$ & $1{,}139$ & $+1{,}70$ & --- \\
Base zero-shot senza Trie & $10{,}09$ & $12{,}26$ & $0{,}823$ & $-2{,}17$ & --- \\
\bottomrule
\end{tabular}
\end{center}

La lettura contraddice l'ipotesi.

\begin{itemize}
  \item Il GRPO dal modello base produce uscite di lunghezza praticamente
        corretta in media (rapporto $1{,}011$, errore firmato $+0{,}14$). Nessuna
        deriva verso la verbosità.
  \item La cella SFT seguito da GRPO produce uscite \emph{più corte} del
        riferimento (rapporto $0{,}886$, errore firmato $-1{,}40$). Se ci fosse
        una pressione della reward sulla lunghezza, agirebbe nella direzione
        opposta a quella prevista.
  \item L'errore \emph{assoluto} è la colonna informativa. L'SFT ha $0{,}09$:
        azzecca la lunghezza quasi sempre. Il GRPO da base ha $2{,}81$ e la
        pipeline SFT$\to$GRPO ha $3{,}59$, con errori firmati piccoli. Cioè: le
        celle con RL sbagliano la lunghezza di circa tre token per volta, in
        entrambe le direzioni, e gli errori si compensano nella media.
\end{itemize}

La conclusione è che il problema delle celle con RL non è un bias di lunghezza
ma una \emph{varianza} di lunghezza. È coerente con l'analisi analitica della
componente \texttt{edit\_validity}, che perde circa $0{,}13$ per token appeso
(equazione~\eqref{eq:append-k}): la reward penalizza attivamente
l'allungamento, quindi la verbosità non può essere il canale di guadagno.

Questa è una ritrattazione dichiarata: l'ipotesi ``la reward premia la
verbosità'' era ragionevole a priori, è stata testata due volte
(analiticamente e sui dati di lunghezza), e va abbandonata.

\section{Osservazione 5: il presunto crollo della pipeline SFT$\to$GRPO}
\label{sec:osservazione5}

Il risultato più scomodo della tabella:

\begin{center}
\begin{tabular}{lrr}
\toprule
Cella & ROUGE-L & EM \\
\midrule
SFT soltanto & $0{,}9749$ & $0{,}8117$ \\
SFT seguito da GRPO & $0{,}5114$ & $0{,}0203$ \\
\bottomrule
\end{tabular}
\end{center}

Applicare il GRPO a una policy già addestrata la \textbf{distrugge}: ROUGE-L da
$0{,}9749$ a $0{,}5114$, exact match da $0{,}8117$ a $0{,}0203$. Non un
peggioramento marginale, un crollo di quaranta punti di ROUGE-L e di quasi
ottanta punti di exact match.

Il meccanismo è quello documentato nella tabella del supporto dei rollout
(sezione~\ref{sec:supporto}). Sulla policy SFT, l'$84{,}8\%$ dei gruppi ha
varianza nulla: reward media $+0{,}9462$, nessun gruppo al valore minimo, tutti
gli otto candidati corretti e con la stessa reward. Su quei gruppi il
vantaggio~\eqref{eq:grpo-adv} è esattamente zero e non c'è aggiornamento utile.

Sul restante $15{,}2\%$, invece, c'è varianza, e l'aggiornamento avviene. Ma
quel $15\%$ è per costruzione il sottoinsieme in cui la policy SFT era
\emph{incerta}, cioè i casi difficili e potenzialmente atipici. Ottimizzare la
reward su un campione così selezionato, con la penalità KL come unico freno
($\beta = 0{,}04$), sposta la policy in una regione che soddisfa meglio la reward
e peggio il riferimento. In assenza di dynamic sampling
(sezione sui limiti di TRL nel Capitolo~\ref{cap:grpo}) non esiste nessun
meccanismo che compensi questa selezione.

Il fenomeno è documentato in letteratura come sovraottimizzazione del modello di
reward~\cite{gao2022overoptimization}: la reward è un surrogato imperfetto, e
oltre un certo punto migliorarla peggiora l'obiettivo vero. La particolarità qui
è che si dispone della misura che lo prevedeva \emph{prima} dell'esperimento:
l'$84{,}8\%$ di varianza nulla era osservabile con una singola sonda sui rollout,
senza addestrare.

\paragraph{Risoluzione: il crollo era in gran parte il guasto~6.} Durante la
stesura di questa relazione è stato individuato e corretto un difetto nella
valutazione delle celle \texttt{sft-grpo} (guasto~6,
sezione~\ref{sec:bug-sft-merge-eval}): l'adattatore della fase SFT non veniva
incluso nel modello ricostruito per l'eval, che quindi valutava
\texttt{base + adattatore GRPO} invece di \texttt{base + SFT + adattatore
GRPO}. I numeri $0{,}9749$/$0{,}5114$ sopra sono stati prodotti con il codice
\emph{non corretto}.

Una campagna successiva, con la correzione in codice fin dall'addestramento
(dettagli in sezione~\ref{sec:sft-merge-non-rimisurato}), dà un esito netto:
la cella \texttt{sft-grpo/few-shot} composta correttamente ottiene ROUGE-L
$0{,}9628$ ed exact match $0{,}7296$ (few-shot nativo), contro $0{,}9681$ e
$0{,}7662$ dell'SFT da sola nella stessa campagna (zero-shot nativo). I due
numeri vengono da modalità di prompting diverse --- la stessa cella SFT
valutata in few-shot, nella stessa campagna, dà $0{,}9191$/$0{,}538$ ---
quindi il confronto non è appaiato per prompting; non è nemmeno il confronto
più favorevole alla cella \texttt{sft-grpo}, che con qualunque delle due
letture resta comunque entro un punto o due di ROUGE-L dal soffitto SFT, non
ai quaranta punti misurati sopra.

\begin{center}
\begin{tabular}{lrr}
\toprule
Cella (campagna con guasto~6, $2\,000$ passi) & ROUGE-L & EM \\
\midrule
SFT soltanto & $0{,}9749$ & $0{,}8117$ \\
SFT seguito da GRPO (\emph{con guasto~6}) & $0{,}5114$ & $0{,}0203$ \\
\midrule
\multicolumn{3}{l}{\emph{Campagna successiva, guasto~6 corretto ($5\,000$ passi, non direttamente comparabile alla precedente):}} \\
SFT soltanto & $0{,}9681$ & $0{,}7662$ \\
SFT seguito da GRPO (\emph{corretto}) & $0{,}9628$ & $0{,}7296$ \\
\bottomrule
\end{tabular}
\end{center}

Il meccanismo dei gruppi a varianza nulla (l'$84{,}8\%$ senza gradiente
utile) resta un fatto reale, misurato direttamente sui rollout indipendentemente
da come il checkpoint viene poi ricomposto per l'eval; ma la sua portata va
ridimensionata. Non è la causa di un crollo di quaranta punti --- quel numero
era in gran parte il guasto~6 --- bensì di una deriva piccola, dell'ordine di
un punto di ROUGE-L, coerente con un aggiornamento che tocca solo il $15{,}2\%$
dei gruppi e resta ancorato dalla penalità KL ($\beta = 0{,}04$). Questa è
un'altra ritrattazione dichiarata, nello stesso spirito
dell'Osservazione~4: l'entità dell'effetto, non il meccanismo che lo produce,
era sbagliata.

\section{Osservazione 6: gli obiettivi ausiliari}
\label{sec:osservazione-ausiliari}

I due obiettivi del Capitolo~\ref{cap:ausiliari} sono stati addestrati come
varianti della cella SFT zero-shot, con un solo fattore cambiato: il termine
ausiliario, di peso $\lambda = 0{,}1$ con rampa lineare sui primi $200$ passi.
La valutazione è la stessa della cella SFT: $2\,000$ prompt del test set,
decodifica vincolata dal Trie, campionamento a $T = 0{,}7$ con $5$ completion per
prompt.

\begin{center}
\begin{tabular}{lrrrr}
\toprule
Cella & Passi (epoche) & ROUGE-L & Exact match & chrF \\
\midrule
SFT soltanto & $10\,000$ ($2{,}24$) & $0{,}9681$ & $0{,}7662$ & $96{,}36$ \\
Massa ammessa & $13\,000$ ($2{,}91$) & $\mathbf{0{,}9769}$ & $\mathbf{0{,}8341}$ & $97{,}55$ \\
Loss strutturata, grafo vero & $5\,600$ ($1{,}25$) & $0{,}9430$ & $0{,}6053$ & $93{,}16$ \\
Controllo, prima versione$^{*}$ & $10\,000$ ($2{,}24$) & $0{,}9693$ & $0{,}7737$ & $96{,}53$ \\
Controllo, versione corretta & \multicolumn{4}{c}{in corso} \\
\bottomrule
\end{tabular}
\end{center}
{\small $^{*}$ Non valido come controllo: si veda sotto e la
sezione~\ref{sec:controllo-mescolato}.}

\paragraph{La durata dell'addestramento non è la stessa, per scelta.} Ogni cella
si ferma per early stopping sulla loss di validazione interna (valutata ogni $200$
passi, pazienza di tre valutazioni), e quella loss è la loss \emph{totale}: LM più
$\lambda$ volte il termine ausiliario. Ciascuna cella si ferma quindi quando smette
di migliorare il \emph{proprio} obiettivo, e le durate differiscono. È una scelta
deliberata e non un difetto del confronto: il confronto misura ogni obiettivo
nelle condizioni in cui il suo stesso criterio di arresto lo porta. Va però
tenuta presente in ogni lettura dei delta qui sotto.

\subsection{La massa ammessa}

Rispetto all'SFT la cella guadagna $+0{,}0088$ di ROUGE-L e $+0{,}0679$ di exact
match. Le due metriche vanno lette in modo diverso:
\begin{itemize}
  \item su \textbf{ROUGE-L} il guadagno è sotto la soglia di circa $0{,}02$ oltre
        la quale una differenza fra celle è distinguibile dalla variabilità di
        addestramento (sezione~\ref{sec:dispersione-run}). Non è un risultato;
  \item su \textbf{exact match} il guadagno è circa $2{,}2$ volte la dispersione
        fra due run della cella SFT ($0{,}0304$). È sopra il rumore osservato, ma
        non di molto, e con un solo seed per cella.
\end{itemize}

Il risultato sull'exact match \emph{contraddice} la previsione dichiarata in
anticipo nel Capitolo~\ref{cap:ausiliari}: la loss è piatta sull'insieme ammesso,
non esprime preferenze fra token legali e quindi, da sola, non dovrebbe migliorare
l'exact match. Prima di concludere che la previsione era sbagliata vanno escluse
due spiegazioni alternative:
\begin{enumerate}
  \item \textbf{più addestramento}: la cella ha fatto il $30\%$ di passi in più
        dell'SFT ($13\,000$ contro $10\,000$), perché il termine ausiliario ha
        ritardato l'early stopping;
  \item \textbf{una distribuzione più affilata sotto campionamento}. Il valore
        della loss è piatto rispetto a ridistribuzioni interne all'insieme
        ammesso, ma il suo gradiente non lo è: per un token ammesso,
        $\partial \mathcal{L}_{\text{AM}} / \partial z_i
        = -\,p_i (1 - P_{\allowed}) / P_{\allowed}$, dove $P_{\allowed}$ è la
        massa dell'insieme. Ogni logit ammesso sale quindi in proporzione alla
        propria probabilità: i token già probabili guadagnano più degli altri,
        e la distribuzione \emph{dentro} l'insieme si affila anche dopo la
        rinormalizzazione del Trie. L'ordinamento dei token ammessi però non
        cambia, perché $z_i - z_j$ si muove nella direzione di $p_i - p_j$.
        Sotto campionamento a $T = 0{,}7$ un modello più affilato si allontana
        meno spesso dalla sua prima scelta, e l'exact match sale anche con la
        stessa prima scelta. Sarebbe un effetto di calibrazione, compatibile
        con la previsione del Capitolo~\ref{cap:ausiliari}.
\end{enumerate}
La seconda spiegazione si verifica valutando le due celle in modo
\emph{deterministico} (greedy: si prende sempre il token più probabile, nessuna
casualità). Poiché l'affilatura conserva l'ordinamento, in greedy il solo termine
ausiliario non dovrebbe cambiare le risposte: se il vantaggio resta, la cella ha
imparato risposte migliori (per effetto indiretto sull'addestramento o per la
durata maggiore); se sparisce, era calibrazione. La valutazione greedy di
entrambe le celle è \textbf{in corso} (celle \texttt{sft/zero-shot-greedy} e
\texttt{ablations/\allowbreak objectives/\allowbreak sft-allowed-mass-greedy}).
Il criterio~1 nella sua forma originale, la sonda in inferenza sulla massa
rimossa dal vincolo, non è ancora stato calcolato.

\subsection{La loss strutturata}

Con il grafo vero la cella è \emph{peggiore} dell'SFT su tutte le metriche:
$-0{,}0251$ di ROUGE-L, oltre la soglia di distinguibilità, e $-0{,}1609$ di
exact match. Si è anche fermata molto prima ($5\,600$ passi, $1{,}25$ epoche),
per la ragione descritta sopra: il termine strutturato pesa circa $0{,}32$ su una
loss di validazione totale di $0{,}40$ (NLL strutturata attorno a $3{,}2$, valore
diagnostico dell'ultimo batch valutato), per cui l'arresto è governato dal
plateau della testa strutturata e non da quello del modello linguistico.

Con i dati disponibili non si possono separare le due letture: modello meno
addestrato, oppure termine strutturato che interferisce con l'apprendimento
della parte linguistica. Un confronto a pari numero di passi non è ricostruibile
dai log, perché la loss LM registrata durante la validazione è quella dell'ultimo
batch e non la media sull'insieme di validazione. In entrambe le letture, in
questa configurazione il risultato non fornisce alcuna evidenza a favore
dell'obiettivo, ed è coerente con il margine modesto stimato in anticipo
($+0{,}0444$ bit, sezione~\ref{sec:qualificazione-strutturale}).

Una prima esecuzione di questa cella e del suo controllo si era fermata a $800$
passi per un guasto, documentato nel Capitolo~\ref{cap:infrastruttura}: la loss
di validazione valeva $+\infty$ a ogni valutazione. I numeri della tabella sono
quelli successivi alla correzione.

\subsection{Il controllo: la prima versione non era un controllo}

La prima versione del controllo mescolato permutava le destinazioni degli archi
del grafo, e così facendo rendeva non rappresentabili i percorsi gold reali
(sezione~\ref{sec:controllo-mescolato}): il termine strutturato è rimasto spento
nel $98\%$ dei passi. La cella si è quindi comportata come un secondo SFT, e i
numeri lo confermano: si ferma allo stesso passo dell'SFT ($10\,000$) e ne
riproduce le metriche entro $0{,}0012$ di ROUGE-L e $0{,}0075$ di exact match,
valori coerenti con la dispersione fra run della sezione~\ref{sec:dispersione-run}.
Come replica dell'SFT è un dato utile; come controllo è inservibile.

Il confronto ``grafo vero contro grafo mescolato'' che la prima versione
sembrava fornire (il grafo mescolato ``vinceva'') misurava in realtà ``termine
strutturato contro nessun termine''. Il verdetto sul criterio~2 dipende dal
controllo corretto, che conserva il supporto del grafo e permuta solo i pesi;
il suo addestramento è \textbf{in corso}.

\section{La decomposizione dell'effetto}
\label{sec:decomposizione}

La domanda che chiude il capitolo: quanto dell'apparente miglioramento è
attribuibile al reinforcement learning?

Si consideri il percorso dal punto di partenza peggiore al miglior risultato con
RL:

\begin{center}
\begin{tabular}{lrl}
\toprule
Passaggio & ROUGE-L & Variazione \\
\midrule
Base zero-shot con Trie & $0{,}1373$ & --- \\
Base few-shot con Trie & $0{,}4664$ & $\mathbf{+0{,}33}$ \quad (prompting, nessun addestramento) \\
GRPO da modello base & $0{,}5998$ & $+0{,}13$ \quad (RL, confuso) \\
\bottomrule
\end{tabular}
\end{center}

Il totale è circa $+0{,}46$, di cui:

\begin{itemize}
  \item \textbf{$0{,}33$ viene dal prompting}, cioè dall'inserimento di tre
        esempi recuperati nel prompt. Nessun parametro aggiornato, nessun tempo
        di GPU per addestramento. Circa il $72\%$ dell'effetto totale.
  \item \textbf{$0{,}13$ è attribuito al RL}, ma in modo \emph{confuso}: la cella
        GRPO da modello base non è isolata dal fattore prompting, e non esiste
        una cella che vari soltanto il metodo di addestramento tenendo fisso il
        prompting.
\end{itemize}

Questa è la confusione di fattori centrale nel disegno originale. Le celle
sperimentali variavano simultaneamente metodo di addestramento e modalità di
prompting, quindi il contributo dei due non è separabile dai dati raccolti.

La conseguenza per la lettura di questa relazione è precisa: \textbf{il fattore
sperimentale più efficace misurato in tutto il progetto è il prompting}, e non
il metodo di addestramento. Se si dovesse riprogettare l'esperimento, il fattore
primario sarebbe la griglia prompting $\times$ decodifica, con il metodo di
addestramento come fattore secondario applicato \emph{dopo} aver identificato il
substrato con supporto migliore (che, secondo la
sezione~\ref{sec:supporto}, è few-shot con Trie: gruppi morti allo $0{,}85\%$).

\section{La dispersione fra run e i limiti degli intervalli di confidenza}
\label{sec:dispersione-run}

Otto celle sperimentali sono state eseguite \emph{due volte}, in occasioni
distinte. Il confronto fra le coppie di run permette una verifica che gli
intervalli di confidenza riportati altrove in questa relazione non possono
fornire, e il risultato impone una precisazione.

\begin{center}
\begin{tabular}{lrrr}
\toprule
Cella & Run 1 & Run 2 & Dispersione \\
\midrule
Ablazione senza grammatica & $0{,}5878$ & $0{,}5717$ & $\mathbf{0{,}0161}$ \\
SFT seguito da GRPO & $0{,}5014$ & $0{,}5114$ & $\mathbf{0{,}0099}$ \\
GRPO da modello base & $0{,}6076$ & $0{,}5998$ & $\mathbf{0{,}0077}$ \\
Ablazione Viterbi & $0{,}5105$ & $0{,}5171$ & $0{,}0066$ \\
Ablazione soft-Viterbi & $0{,}5048$ & $0{,}5090$ & $0{,}0041$ \\
SFT soltanto & $0{,}9785$ & $0{,}9749$ & $0{,}0036$ \\
Ablazione stack storico & $0{,}5128$ & $0{,}5101$ & $0{,}0027$ \\
Ablazione struttura & $0{,}5108$ & $0{,}5101$ & $0{,}0007$ \\
\bottomrule
\end{tabular}
\end{center}

Il termine di confronto è l'ampiezza tipica degli intervalli di confidenza
ottenuti per bootstrap, che su ROUGE-L è di circa $\pm 0{,}0044$
(l'intervallo misurato sulla cella senza vincolo è
$[0{,}3605,\; 0{,}3692]$).

\textbf{Tre celle su otto hanno una dispersione fra run superiore all'intera
ampiezza dell'intervallo di confidenza.} Sull'ablazione senza grammatica la
dispersione ($0{,}0161$) è quasi quattro volte l'intervallo.

\subsection{Perché accade, e che cosa implica}

Non è un'anomalia: è una conseguenza di che cosa il bootstrap misura. Il
bootstrap clusterizzato per prompt descritto nel Capitolo~\ref{cap:metriche}
ricampiona i \emph{prompt di valutazione}, quindi stima l'incertezza dovuta alla
scelta del campione di test. Non stima, e non può stimare, la variabilità
dovuta all'addestramento: seed, ordine di presentazione dei dati, non
determinismo delle operazioni su GPU, rumore di quantizzazione.

Segue una regola di lettura che si applica a tutta la relazione:

\begin{quote}
Gli intervalli di confidenza riportati sono \textbf{limiti inferiori}
dell'incertezza reale. Una differenza fra celle inferiore a circa $0{,}02$ di
ROUGE-L non è distinguibile dalla variabilità di addestramento con i dati
disponibili.
\end{quote}

\subsection{La conseguenza sul confronto SFT contro regola}

Questa osservazione tocca direttamente l'affermazione più delicata della
relazione. Il Capitolo~\ref{cap:sft} riporta:

\begin{center}
\begin{tabular}{lcc}
\toprule
Delta SFT $-$ regola & Valore & Intervallo bootstrap $95\%$ \\
\midrule
ROUGE-L & $+0{,}0064$ & $[+0{,}0038,\; +0{,}0098]$ \\
Exact match & $+0{,}2252$ & $[+0{,}2025,\; +0{,}2510]$ \\
\bottomrule
\end{tabular}
\end{center}

Ma la cella SFT, da sola, varia di $0{,}0036$ fra i suoi due run su ROUGE-L, e
di $\mathbf{0{,}0304}$ su exact match ($0{,}8421$ contro $0{,}8117$). Ne segue
una lettura differenziata:

\begin{itemize}
  \item il delta su \textbf{ROUGE-L} ($+0{,}0064$) è dello stesso ordine di
        grandezza della dispersione propria dell'SFT ($0{,}0036$). Va
        \textbf{declassato}: non è un vantaggio stabilito, è un vantaggio
        compatibile con il rumore di addestramento;
  \item il delta su \textbf{exact match} ($+0{,}2252$) resta ampiamente
        robusto: è circa sette volte la dispersione osservata ($0{,}0304$).
        L'affermazione che l'SFT abbatta il tasso di errore dal $41\%$ al
        $19\%$ regge.
\end{itemize}

Questo non indebolisce la tesi centrale della relazione: la rafforza. Se su un
corpus il vantaggio di un modello addestrato su $73$ mila esempi rispetto a una
tabella di corrispondenze non è distinguibile dal rumore quando misurato con la
metrica di riferimento della letteratura, allora quella metrica ha esaurito la
propria capacità discriminante in modo ancora più netto di quanto la
sezione~\ref{sec:decomposizione} suggerisse.

\subsection{Che cosa si sarebbe dovuto fare}

La procedura corretta, e la raccomandazione per chi riprenda il lavoro, è di
eseguire ogni cella con almeno tre seed distinti e riportare media e deviazione
standard fra i seed, invece di un singolo run con l'intervallo bootstrap. Il
costo è triplo, ma è l'unico modo di distinguere un effetto reale da una
fluttuazione di addestramento.

Questa raccomandazione ha \textbf{precedenza} su quella di aumentare il numero
di prompt valutati. Passare da $2\,000$ a $5\,000$ prompt restringe
l'intervallo bootstrap da circa $\pm 0{,}004$ a $\pm 0{,}003$, ma agisce sulla
componente di incertezza \emph{minore}: non tocca la dispersione da
addestramento, che su tre celle su otto è già superiore all'intervallo intero.
Un campione di valutazione più ampio produce quindi intervalli più stretti
attorno a un valore che resta instabile fra run, cioè una falsa impressione di
precisione. L'ordine corretto degli interventi è prima i seed multipli, poi
eventualmente il campione più ampio.

Le otto coppie di run disponibili sono state prodotte come conseguenza
dell'infrastruttura, non per disegno, ed è per questo che sono soltanto due per
cella. Il fatto che siano sufficienti a mettere in discussione un intervallo di
confidenza pubblicato in questa stessa relazione è la migliore dimostrazione
della loro necessità.

\section{Riepilogo per non-copy accuracy}

Poiché la non-copy-token accuracy è la metrica con maggiore potere discriminante
su questo corpus (Capitolo~\ref{cap:metriche}), vale la pena chiudere con
l'ordinamento che essa induce:

\begin{center}
\begin{tabular}{lr}
\toprule
Sistema & Non-copy accuracy \\
\midrule
SFT soltanto & $0{,}9660$ \\
Baseline a regole & $0{,}9424$ \\
GRPO da modello base & $0{,}4641$ \\
Base zero-shot con Trie & $0{,}0432$ \\
Base zero-shot senza Trie & $0{,}0006$ \\
\bottomrule
\end{tabular}
\end{center}

L'ordinamento è lo stesso che si ottiene con ROUGE-L, con un'eccezione
significativa: le prime due righe si stringono ($0{,}9660$ contro $0{,}9424$,
cioè $+0{,}0236$) e le ultime due si separano in modo interpretabile. In
particolare, il GRPO da modello base risolve correttamente meno della metà delle
posizioni non banali, contro il $94\%$ di una tabella di corrispondenze.

Nessuna delle celle con reinforcement learning si avvicina alla regola. È il
risultato del progetto, ed è negativo.
